\begin{textsample}{POS Dim 1 – vem\_ed\_17 – Score 6.00 – vem\_ed\_17/t077.txt}  \label{ex:f1_pos_135}
PERCEPÇÕES DOS ALUNOS: PESQUISA 2º SEMESTRE 2022 Desde o segundo semestre de 2020, a equipe dos Projetos Colaborativos Internacionais (PCIs / Cesu) realiza pesquisas de percepção com alunos e professores participantes de PCIs, para avaliação e melhoria \textbf{contínua} dos \textbf{processos} de gestão dos Intercâmbios Virtuais.
Os resultados mostram a manutenção das boas percepções sobre os projetos ao \textbf{longo} do tempo, tanto por discentes \textbf{quanto} por docentes envolvidos nas colaborações internacionais.
A pesquisa de percepção, enviada para os 1.683 participantes de PCIs entre agosto e dezembro de 2022, obteve 632 respostas de estudantes – dos quais 69 não autorizaram a divulgação dos dados e tiveram suas respostas eliminadas para fins de publicações.
Para a maioria dos respondentes (82 \%), este \textbf{foi} o primeiro PCI; 13 \% afirmaram \textbf{que} \textbf{foi} o segundo.
Dentre os principais desafios, os dois mais mencionados \textbf{foram} gerenciamento de tempo (248 respostas) e comunicação (161), repetindo a tendência dos anos \textbf{anteriores}.
Nesta página, confira a síntese dos resultados relativos aos discentes.
Na página 7, conheça a percepção dos docentes.
Pesquisas \textbf{anteriores} (segundo semestre de 2020, primeiro semestre de 2021, segundo semestre de 2021 e primeiro semestre de 2022) \textbf{foram} relatadas nas edições 5, 8, 9, 12 e 14 de VEm.
Confira em: https://cesu.cps.sp.gov.br/vem-newsletter-bimestral-dos-projetos-colaborativos-internacionais-pci/ 632 respondentes 82 \% 1ª vez participando Alunos \textbf{que} recomendariam aos colegas participar de um PCI 93 \% 86 \% afirmam \textbf{que} a “experiência com o PCI aumentou meu interesse em permanecer no curso até a conclusão”. 87 \% \textbf{reconhecem} \textbf{que} sua \textbf{competência} no \textbf{idioma} estrangeiro melhorou. 85 \% \textbf{reconhecem} \textbf{que} “a experiência com o PCI aumenta as chances de sucesso no \textbf{mercado} de trabalho”. 87 \% avaliam \textbf{que} a \textbf{interação} com os colegas do grupo brasileiro \textbf{foi} ótima ou boa 67 \% a \textbf{interação} com seus colegas do grupo estrangeiro \textbf{foi} ótima ou boa.
Realizaram PCIs em português 6 \% Realizaram PCIs em espanhol 26 \% Realizaram PCIs em inglês 68 \%

% matched lemmas: anterior, competência, contínuo, idioma, interação, longo, mercado, processo, quanto, que, reconhecer, ser
\end{textsample}
